\subsubsection{sự phản xạ sóng}
\begin{itemize}
\item Với đầu phản xạ cố định: sóng tới và sóng phản xạ \textit{ngược pha} tại điểm phản xạ.
\item Với đầu phản xạ tự do: sóng tới và sóng phản xạ \textit{cùng pha} tại điểm phản xạ.
\end{itemize}
\subsubsection{Định nghĩa}
Sóng có các nút sóng và bụng sóng cố định trong không gian gọi là \textbf{sóng dừng}.
\begin{itemize}
\item Những điểm luôn luôn đứng yên gọi là \textbf{\textit{nút}}.
\item Những điểm luôn luôn dao động với biên độ lớn nhất là \textbf{\textit{bụng}}.
\end{itemize}
\Note{Sóng dừng thường là giao thoa giữa sóng tới và sóng phản xạ}
\subsubsection{Điều kiện để có sóng dừng trên dây}
\paragraph{Vị trí các nút và bụng}
immini[thm]{\begin{itemize}
\item Khoảng cách giữa hai nút sóng hoặc hai bụng sóng bất kì:
\begin{align*}
d_{bb}=d_{nn}=k.\dfrac{\lambda}{2}; \quad k=1,2,3,\ldots
\end{align*}
\lefthand\ Khoảng cách giữa hai nút sóng hoặc hai bụng sóng liền kề là $\dfrac{\lambda}{2}$.
\item Khoảng cách giữa một nút sóng với một bụng sóng bất kì
\begin{align*}
d_{nb}=(2k+1).\dfrac{\lambda}{4};\quad k=0,1,2,3,\ldots
\end{align*}
\lefthand\ Khoảng cách giữa một nút sóng với một bụng sóng liền kề là $\dfrac{\lambda}{4}$
\end{itemize} }{\begin{tikzpicture}[scale=1,thick,>=stealth']
\fill [pattern = north east lines] (8,-0.4) rectangle (8.2,0.4);
\fill [pattern = north east lines] (8,-2.2) rectangle (8.2,-1.3);
\draw[black,line width=1pt] (1,0.5) cos (2,0) sin (3,-0.5) cos (4,0) sin(5,0.5) cos(6,0) sin(7,-0.5) cos(8,0);
\draw[black,line width=0.4pt] (1,0.4) cos (2,0) sin (3,-0.4) cos (4,0) sin(5,0.4) cos(6,0) sin(7,-0.4) cos(8,0);
\draw[black,line width=0.4pt] (1,0.3) cos (2,0) sin (3,-0.3) cos (4,0) sin(5,0.3) cos(6,0) sin(7,-0.3) cos(8,0);
\draw[black,line width=0.4pt] (1,0.2) cos (2,0) sin (3,-0.2) cos (4,0) sin(5,0.2) cos(6,0) sin(7,-0.2) cos(8,0);
\draw[black,line width=0.4pt] (1,0.1) cos (2,0) sin (3,-0.1) cos (4,0) sin(5,0.1) cos(6,0) sin(7,-0.1) cos(8,0);
\draw[black,dashed,line width=1pt] (1,-0.5) cos (2,0) sin (3,0.5) cos (4,0) sin(5,-0.5) cos(6,0) sin(7,0.5) cos(8,0) ;
\draw[black,dashed,line width=0.4pt] (1,-0.4) cos (2,0) sin (3,0.4) cos (4,0) sin(5,-0.4) cos(6,0) sin(7,0.4) cos(8,0) ;
\draw[black,dashed,line width=0.4pt] (1,-0.3) cos (2,0) sin (3,0.3) cos (4,0) sin(5,-0.3) cos(6,0) sin(7,0.3) cos(8,0) ;
\draw[black,dashed,line width=0.4pt] (1,-0.2) cos (2,0) sin (3,0.2) cos (4,0) sin(5,-0.2) cos(6,0) sin(7,0.2) cos(8,0) ;
\draw[black,dashed,line width=0.4pt] (1,-0.1) cos (2,0) sin (3,0.1) cos (4,0) sin(5,-0.1) cos(6,0) sin(7,0.1) cos(8,0) ;
\fill[black] (2,0) circle (1pt); \fill[black] (4,0) circle (1pt);\fill[black] (6,0) circle (1pt);\fill[black] (8,0) circle (1pt);
\draw[dashed] (1,0)--(8,0);
\draw[dashed,line width=0.5pt](4,1)--(4,0);\draw[dashed,line width=0.5pt](6,1)--(6,0);\draw[dashed,line width=0.5pt](7,1)--(7,0);\draw[<->](4,0.7)--(6,0.7);\draw[<->](6,0.7)--(7,0.7);\node [above] at (5,0.7) {$\dfrac{\lambda}{2}$}; \node [above] at (6.5,0.7) {$\dfrac{\lambda}{4}$};
\draw[line width=1.5pt](1,-1.8)--(8,-1.8); \fill[black] (1,-1.8) circle (1pt); \fill[black] (4.5,-1.8) circle (1pt); \draw[dashed] (4.5,-1.8)--(4.5,-1.1); \draw[<->](4.5,-1.6)--(8,-1.6);\node [above] at (6.25,-1.6) {d}; \node [below] at (1,-1.8) {A}; \node [below] at (4.5,-1.8) {M}; \node [right] at (8.15,-1.8) {B};
\draw[line width=1.5pt](8,-0.4)--(8,0.4);
\draw[line width=1.5pt] (8,-2.2)--(8,-1.3);
\draw[->,line width=1pt](2,0.3)--(2,0);\node [above] at (2,0.3) {\normalsize{nút sóng}};
\draw[->,line width=1pt](3,0.8)--(3,0.5);\node [above] at (3,0.8) {\normalsize{bụng sóng}};
\draw[->,line width=1pt](4.2,-0.6)--(4.8,-0.2);\node [below] at (4.2,-0.6) {\normalsize{bó sóng}};
\draw[->,black,line width=1pt] (1,-1.4) cos (1.1,-1.5) sin (1.2,-1.6) cos (1.3,-1.5) sin(1.4,-1.4) cos(1.5,-1.5) sin(1.6,-1.6) cos(1.7,-1.5)--(2.2,-1.5);
\node [above] at (2,-1.5) {\normalsize{Sóng tới}};
\draw[->,black,line width=1pt] (8,-2.2) cos (7.9,-2.1) sin (7.8,-2) cos (7.7,-2.1) sin(7.6,-2.2) cos(7.5,-2.1) sin(7.4,-2) cos(7.3,-2.1) sin (7.2,-2.2) cos (7.1,-2.1)--(6.8,-2.1); \node [below] at (6.7,-2.1) {\normalsize{Sóng phản xạ}};
\end{tikzpicture}}
\paragraph{Điều kiện sóng dừng}
\begin{itemize}
\item Dây cố định hai đầu: \bth{\boder{\ell=k.\dfrac{\lambda}{2}}} với $k = 1,\ 2,\ 3...\Rightarrow\begin{cases}
\text{số bụng}=k\\
\text{số nút}=k+1\\
\end{cases}$
\notebox{\textit{Để có sóng dừng trên dây hai đầu cố định thì chiều dài dây phải bằng số nguyên lần nửa bước sóng}}
\item Một đầu dây cố định, một đầu tự do: \bth{\boder{\ell=(2k+1).\dfrac{\lambda}{4}}} với $k = 0,\ 1,\ 2,\ 3....\Rightarrow\begin{cases}
\text{số bụng}=k+1\\
\text{số nút}=k+1\\
\end{cases}$
\notebox{\textit{Để có sóng dừng trên dây một đầu cố định, một đầu tự do thì chiều dài dây phải bằng số lẻ lần một phần tư bước sóng}}
\end{itemize}